LLM-based code-review agents are increasingly entrusted with gating code
changes, yet current benchmarks largely measure whether an agent produces the
correct accept/reject label without separating two distinct questions: whether
the \emph{evidence} presented is intact and trustworthy, and whether it is
\emph{sufficient} to justify a decision at all. We introduce
\textbf{GroundedReview}, a benchmark that (i) casts code review as a three-way
decision---\ACCEPT, \REJECT, or \NEED---and (ii) evaluates the same real
issue--patch pair under several evidence conditions that independently control
which test outcomes are visible, whether the agent may actively fetch
artifacts, and whether provenance and freshness are verifiable. We distinguish
\emph{Oracle Gold} (the decision warranted by a fresh, scoped pytest run) from
\emph{Condition Gold} (the decision warranted by the evidence actually
granted, which is legitimately \NEED\ on restricted surfaces). On ten SymPy
issue--patch pairs served to two models (\texttt{deepseek-flash} and
\texttt{gpt-4.1-mini}) across five conditions, a corrective replication shows
that repairing protocol-construct defects in our own instrument moved one
model's sufficient-evidence abstention from $9/10$ to $1/10$---an
associational, not causal, contrast---whereas the other model, under a
restricted ``strict'' surface whose correct answer is \NEED, returned
unsupported \ACCEPT\ on $7/10$ issues, upgrading non-decisive evidence into
acceptance. A traceable case shows an agent accepting a patch after an
already-visible fail-to-pass test had failed. We further contribute eight
newly constructed, independently de-duplicated issues with 40 per-condition
gold certificates and a non-interference-checked evidence projection, and we
report a \emph{capped, exploratory} 80-cell run on these independent issues
(two models $\times$ five conditions, spend \$0.258 under \$2.00/\$1.00
hard caps, zero retries or errors). On this second, disjoint issue set the
restricted-evidence pattern recurs: one model abstains on all $8/8$ strict
cells, while the other upgrades the non-decisive strict surface into
unsupported \ACCEPT\ on $6/8$ issues (CP95 $[0.35,0.97]$), including both
oracle-\REJECT\ issues; under sufficient tier-A evidence both models match
the oracle on $8/8$ issues. Because this run uses a different issue set and
time, the recurrence is an associational cross-set replication, not a causal
estimate. All labels remain \textbf{PROVISIONAL} pending independent human
arbitration (\texttt{annotator2 = NOT\_OBTAINED}); we therefore present a
measurement framework and clearly bounded evidence rather than
population-level claims about model behaviour.
